% from main.tex
%%%%%%%%%%  scan idx 506  |  folio 493  %%%%%%%%%%
\textbf{Proposition 5.2.} On a sur \(P\) la suite exacte :
\[
0\longrightarrow\Omega^1_{P/Y}\longrightarrow f^*(E)(-1)\longrightarrow\mathcal O_P\longrightarrow0,
\]
i.e. on a un isomorphisme
\[
\Omega^1_{P/Y}\simeq F(-1),\quad\text{où}\quad F=\Ker\bigl(f^*(E)\longrightarrow\mathcal O_P(1)\bigr).
\]

On considère le fibré vectoriel sur \(Y\) défini par \(E\), et on note \(Z\) le complémentaire de la section nulle de ce fibré. On a le diagramme commutatif :
\[
\begin{tikzcd}[column sep=large,row sep=large]
P \arrow[dr,"f"'] && Z \arrow[ll,"h"'] \arrow[dl,"g"]\\
& Y &
\end{tikzcd}
\]
Le morphisme \(h\) identifie \(Z\) au \(P\)-schéma affine défini par le faisceau d'algèbres \(B=\bigoplus_{n\in\mathbb Z}\mathcal O_P(n)\) ; \(B\) est une \(\mathcal O_P\)-Algèbre graduée.

Comme \(f,g,h\) sont des morphismes lisses, on peut écrire la suite exacte des faisceaux de différentielles relatives ; comme \(Z\) est affine sur \(P\), on peut écrire cette suite en termes de \(B\)-Modules; on obtient ainsi :
\begin{equation}\tag{5.2.1}
0\longrightarrow
 \Omega^1_{P/Y}\otimes_{\mathcal O_P}B
\longrightarrow
 \Omega^1_{B/Y}
\longrightarrow
 \Omega^1_{B/\mathcal O_P}
\longrightarrow0 .
\end{equation}
Or cette suite exacte est une suite exacte graduée de \(B\)-Modules gradués. En effet, elle résulte du diagramme commutatif :
\[
\begin{tikzcd}[column sep=large,row sep=large]
S(E) \arrow[rr] && \displaystyle\bigoplus_{n\in\mathbb Z}S(E)(n)\simeq S(E)[T,T^{-1}]\\
& \mathcal O_X \arrow[ul] \arrow[ur] &
\end{tikzcd}
\]
où les flèches sont des homomorphismes d'algèbres, compatibles avec la graduation de \(S(E)[T,T^{-1}]\) (graduation définie par le degré par rapport à l'indéterminée \(T\)), les deux autres algèbres étant supposées réduites au degré \(0\).

Le
